\scanpage{26}
\uline{Cas particulier}

$A=\operatorname{Sym}_{\Lambda}(M)$ où $M$ est un $\Lambda$-module

\deleted{$\unsure{M}$} $J=\operatorname{Sym}_{\Lambda}^{+}(M)$ l'idéal d'augmentation.

Alors le pb universel est le même que celui d'envoyer $\underline{M}$ dans une $\Lambda$-algèbre à puissances divisées, on trouve
\[
\Gamma(\operatorname{Sym}_{\Lambda}(M),\operatorname{Sym}_{\Lambda}^{+}(M))\simeq\Gamma_{\Lambda}^{\unsure{+|*}}(M)
\]
Ainsi, en tant qu'algèbre \deleted{\unsure{de $\Lambda$-anneau}} à puissances divisées $\Gamma_{\Lambda}(M)$ dépend fonctoriellement de l'anneau \uline{augmenté} $\operatorname{Sym}_{\Lambda}^{\unsure{+|*}}(M)$ \unsure{---}

\deleted{Si $A=\Lambda[T_1,\ldots,T_n]$}

\margin{plus haut au § 1, aux fonctorialités\ldots}

\margin{\unsure{\S\,1}} [Soit $A$ un \deleted{$\Lambda$-algèbre} \inserted{anneau avec idéal $J$} à puissances divisées, considérons des éléments $(x_i)_{i\in I}$ \inserted{de $J$} et proposons-nous de \unsure{trouver} les homs $(A,J)\longrightarrow(A'',J'')$ avec $A'$ un \deleted{$\Lambda$-algèbre} \inserted{anneau à idéal $J'$} à puissances divisées, qui annulent les $x_i$. Existe-t-il un $(A',J')$ universelle? Oui, il suffit de prendre le quotient de $A$ par l'idéal \inserted{$K$} engendré par les $\pi^n(x_i)$ $i\in I$, $n\geq1$, i.e. les expressions
\[
\sum \deleted{(\unsure{\lambda_i})}\deleted{\unsure{\lambda}}(\lambda_{ni}\,\pi^n(x_i)),\quad\inserted{\text{avec}}\deleted{(\unsure{\lambda_i\pi^n(x_i)})}
\]
\deleted{$\unsure{\lambda_i,\lambda}$} $\lambda_{ni}\in A$. On \unsure{trouve} ainsi le plus petit idéal $K$ stable par les $\pi^n$ $(n\geq1)$,
\mspage{25}
